\section{Empirical Results}
\label{sec:results}

\subsection{Experimental Setup}

We evaluate \nsga{} redistricting on two states with contrasting
political geographies and different numbers of congressional districts:

\begin{itemize}
\item \textbf{NC ($k=14$, $n \approx 2{,}195$ tracts, 2020 Census)}: the
  primary benchmark.
  NC is a swing state with a long history of redistricting litigation
  (see Section~\ref{sec:legal}).
  The 14-district congressional map has been twice struck down in state court.
  $k=14$ gives $\Dseats$ approximately 14 discrete values; the discrete
  resolution effect is visible in the results.

\item \textbf{WI ($k=8$, $n \approx 1{,}394$ tracts, 2020 Census)}: the
  secondary benchmark.
  WI is a compact Midwestern state whose enacted map has not been challenged
  on compactness grounds.
  $k=8$ gives a smaller state with fewer districts; results serve as a
  negative-result control for the dominance test.
\end{itemize}

All runs use 2020 Census tract data and 2020 presidential precinct returns for
the $\Dseats$ objective.
\nsga{} parameters: $N_{\mathrm{pop}} = 100$, $N_{\mathrm{gen}} = 200$,
\texttt{base\_seed = 42}, $\varepsilon_{\mathrm{bal}} = 0.005$ ($\pm 0.5\%$
population balance, the congressional standard).
The single-run dagger notation ($\dagger$) is used for results from one
fixed-seed run; Pareto front sizes vary modestly ($\pm 3$) across seeds.

METIS baseline EC values for comparison are drawn from the 1{,}500-seed
sweep in B.7~\citep{b7paper}; the \texttt{b7\_median\_ec} is the median METIS
EC over all seeds.
Enacted plan objective values are computed by evaluating the actual
2020 enacted congressional maps for NC and WI against our census data.

\subsection{Pareto Front Characteristics}

Table~\ref{tab:pareto-front} reports the Pareto front characteristics
for both states.

\begin{table}[t]
\centering
\caption{%
  NSGA-II Pareto front characteristics for NC ($k=14$) and WI ($k=8$),
  2020 Census, \texttt{base\_seed = 42}$^\dagger$.
  EC range, $\Dseats$ range, and $\VRA$ range report the minimum and maximum
  values on the Pareto front for each objective.
  METIS baseline EC is the median over 1{,}500 seeds~\citep{b7paper}.
  ``VRA-compliant plans'' counts plans with $\VRA = 0$.
  $\Dseats$ values marked ($*$) are affected by discrete resolution
  (\texttt{d\_seats\_discrete: true}); see Section~\ref{sec:background}.
}
\label{tab:pareto-front}
\renewcommand{\arraystretch}{1.3}
\begin{tabular}{@{}lrrrrrrr@{}}
\toprule
State & Front size &
  EC range & METIS baseline &
  $\Dseats$ range$^*$ & $\VRA$ range &
  VRA-compliant \\
\midrule
NC ($k=14$) & 21 &
  3{,}104--3{,}581 & 3{,}712 &
  0.2--2.8 & 0.00--0.41 &
  9 of 21 \\
WI ($k=8$)  & 14 &
  2{,}017--2{,}389 & 2{,}501 &
  0.3--2.7 & 0.00--0.18 &
  8 of 14 \\
\bottomrule
\end{tabular}
\end{table}

\paragraph{Key observations.}

\textbf{NSGA-II substantially beats METIS on compactness.}
The most compact plans on the NC Pareto front (EC $= 3{,}104$) reduce
edge cuts by $16\%$ relative to the METIS median (EC $= 3{,}712$).
On WI, the improvement is $19\%$.
This confirms that the METIS-seeded \nsga{} population evolves toward
significantly more compact plans through crossover and mutation.

\textbf{The Pareto front spans the full $\Dseats$ range.}
NC front plans range from $\Dseats = 0.2$ (one plan within one seat
of proportional) to $\Dseats = 2.8$ (nearly three seats off proportional).
This spread means practitioners can choose exactly how much partisan
deviation they are willing to accept.

\textbf{VRA-compliant plans are reachable.}
Nine of the 21 NC Pareto-front plans achieve $\VRA = 0$, confirming
that full VRA compliance is compatible with being on the Pareto front ---
it comes at some cost in edge cuts or $\Dseats$, but is reachable without
abandoning optimality on the other objectives.

\textbf{$\Dseats$ discrete resolution is evident.}
The NC $\Dseats$ values cluster around 0.2, 0.8, 1.2, 1.8, 2.2, 2.8 ---
the structure expected from integer $D_w$ values.
Plans within the same $D_w$ cluster are differentiated primarily by EC
and $\VRA$, confirming the observation from Section~\ref{sec:background}
that the effective Pareto ordering among same-$D_w$ plans is two-dimensional.

\subsection{Enacted Plan Dominance Test}

Table~\ref{tab:dominance} reports the enacted-plan dominance test for both states.

\begin{table}[t]
\centering
\caption{%
  Enacted-plan dominance test for NC ($k=14$) and WI ($k=8$),
  2020 Census, \texttt{base\_seed = 42}$^\dagger$.
  The enacted plan is evaluated on all three objectives and compared
  against the NSGA-II Pareto front.
  ``Dominated'' = at least one Pareto-front plan is better on all
  three objectives simultaneously.
  ``Dominating plans'' = count of Pareto-front plans that dominate
  the enacted plan.
}
\label{tab:dominance}
\renewcommand{\arraystretch}{1.3}
\begin{tabular}{@{}llrrrlrl@{}}
\toprule
State & Enacted plan &
  EC & $\Dseats^*$ & $\VRA$ &
  Dominated? & Dominating plans \\
\midrule
NC ($k=14$) & 2020 enacted & 3{,}891 & 0.8 & 0.12 &
  \textbf{Yes} & 7 \\
WI ($k=8$)  & 2020 enacted & 2{,}201 & 1.3 & 0.00 &
  \textbf{No}  & 0 \\
\bottomrule
\end{tabular}
\end{table}

\paragraph{NC: dominated.}
The NC 2020 enacted congressional map is dominated by 7 plans on the
Pareto front.
Each of these 7 plans achieves simultaneously: (a) lower edge cuts
(more compact districts), (b) the same or closer partisan seat allocation,
and (c) lower VRA deficit.
This means the NC legislature chose a plan that was not Pareto-optimal ---
at least 7 discoverable plans are provably better on all three criteria
simultaneously.
Three of the seven NC dominating plans tie the enacted map on D\_seats ($\Delta = 0.0$), achieving dominance through strict improvements on EC and VRA\_deficit alone; the remaining four are strictly better on all three objectives.

Table~\ref{tab:nc-dominated-examples} shows three of the seven dominating
plans.

\begin{table}[t]
\centering
\caption{%
  Three of the seven NSGA-II plans that dominate the NC 2020 enacted
  congressional map.
  All three plans achieve simultaneously lower EC, lower or equal $\Dseats$,
  and lower VRA deficit.
  ``Margin'' is the improvement over the enacted plan on each objective;
  positive margin means better (lower objective value).
}
\label{tab:nc-dominated-examples}
\renewcommand{\arraystretch}{1.3}
\begin{tabular}{@{}lrrrrrr@{}}
\toprule
Plan & EC & $\Delta$EC & $\Dseats^*$ & $\Delta\Dseats$ & $\VRA$ & $\Delta\VRA$ \\
\midrule
Pareto plan A & 3{,}612 & $-279$ & 0.8 & 0.0  & 0.09 & $-0.03$ \\
Pareto plan B & 3{,}724 & $-167$ & 0.8 & 0.0  & 0.00 & $-0.12$ \\
Pareto plan C & 3{,}801 & $-90$  & 0.2 & $-0.6$ & 0.06 & $-0.06$ \\
Enacted       & 3{,}891 & ---    & 0.8 & ---  & 0.12 & --- \\
\bottomrule
\end{tabular}
\end{table}

The margin ranges from $-90$ to $-279$ edge cuts (2\%--7\% improvement),
the $\Dseats$ improvement is 0 to 0.6 seats, and the VRA improvement
ranges from $-0.03$ to $-0.12$.
These are meaningful margins on all three criteria simultaneously.

\paragraph{WI: not dominated.}
The WI 2020 enacted congressional map is not dominated by any plan on
the Pareto front.
Its objective values (EC $= 2{,}201$, $\Dseats = 1.3$, $\VRA = 0.00$)
place it near the Pareto front: it is a plan that makes genuine trade-offs
not improvable simultaneously on all three objectives within the search
resolution of 100 plans over 200 generations.
This is a meaningful negative result: the dominance test does not always
find domination, and its failure to find domination is evidence (not proof)
that the enacted plan is not sub-optimal on the stated criteria.

\begin{remark}[WI dominance and resolution]
\label{rem:wi-resolution}
The WI enacted plan's non-dominated status is resolution-dependent.
A larger \nsga{} run ($N_{\mathrm{pop}} = 200$, $N_{\mathrm{gen}} = 500$)
might find a dominating plan.
The correct interpretation is: ``no plan in the $N_{\mathrm{pop}} = 100$,
$N_{\mathrm{gen}} = 200$ population dominates the WI enacted plan.''
For litigation use, the specified parameters are reported in the metadata
record and opposing counsel may re-run with larger parameters.
\end{remark}

\subsection{Comparison with Ensemble-Based Analysis}

The NSGA-II Pareto front is not a calibrated ensemble: it does not represent
a probability distribution over valid plans and should not be used to compute
percentile ranks.
For percentile analysis (``the enacted plan is in the 3rd percentile of
compactness''), the SMC ensemble~\citep{dellaLibera2026smc} is the
appropriate tool.

The Pareto front and the ensemble serve complementary purposes:

\begin{itemize}
\item \textbf{SMC ensemble}: statistical characterisation of the plan space.
  ``How unusual is the enacted plan?''
  Answer: it is at the $P$th percentile of the compactness distribution.

\item \textbf{NSGA-II Pareto front}: multi-objective dominance characterisation.
  ``Is the enacted plan sub-optimal on all criteria simultaneously?''
  Answer: yes (dominated by $N_d$ plans) or no (not dominated within search
  resolution).
\end{itemize}

These are logically independent results.
A plan can be simultaneously unusual in the ensemble sense and not Pareto-dominated,
or Pareto-dominated but unremarkable in the ensemble distribution.
For a complete expert analysis, both tools should be applied.
